
@article{3731171a42664e70b23edeb86a112e68,
title = "Boomgericht bosbeheer en kwaliteitshout",
abstract = "In ieder bos is het mogelijk om kwaliteitshout te telen. De blesser moet vooral de juiste bomen kunnen herkennen en met gerichte ingrepen hun ontwikkeling bijsturen. Ook als de verjonging zich wat spaarzamer aandient en takkiger is, zoals dat onder een kronendak vaker het geval is, streven we naar kwaliteitshout. Hoe komen we met boomgericht bosbeheer vanuit onze huidige bossen in dat structuurrijke en weerbare bos met kwaliteitshout?",
author = "Etienne Thomassen",
year = "2019",
language = "Nederlands",
pages = "7--11",
journal = "Vakblad Natuur Bos Landschap",
issn = "1572-7610",
number = "153",
}
